% ----------------------------------------------
\chapter{Travail Preliminaire}


% ############################################

	
	\textcolor{blue}{Reponse 1: Montrons que $ \Delta P = P_{A}-P{B} = \rho g ( h_{A} - h{B}) + \rho g ( z_{A} - z_{B})   $ }

		
		On applique la loi de l'hydrostatique :
		\[
		\nabla P = \rho \vec{g} \quad \Rightarrow \quad P = -\rho g z + Cste
		\]
		
		Pour les prises de pression :
		\begin{align*}
			P_A &= -\rho g z_A + Cste \\
			P_B &= -\rho g z_B + Cste
		\end{align*}
		
		A la surface :
		\[
		P_0 = -\rho g h_A + Cste \quad d'ou \quad 		Cste = P_0 + \rho g h_A 
		\]
		
		
		Ainsi :
		\begin{align*}
			P_A &= P_0 + \rho g (- z_A + h_A) \\
			P_B &= P_0 + \rho g (- z_B + h_B)
		\end{align*}

		On obtient alors :
		\[
		P_A - P_B = \rho g \left[ (z_A + h_A) - (z_B + h_B) \right] = \rho g (h_A - h_B) + \rho g (z_B - z_A)
		\]
	



	\textcolor{blue}{Reponse 2: Déduisons les valeurs à mesurer pour déterminer expérimentalement $Q$}

		Pour établir la loi expérimentale $Q = f(\Delta P)$, on doit mesurer : 
		
		\begin{enumerate}
			\item 	le debit volumique Q 
			\item	Les hauteurdeau $h_A et h_B$ 
			\item 	la differece d altitude $z_B et z_A$ entre les prises 
		\end{enumerate}
		

\section{Manipulation } 

	Les expériences ont été faites et les valeurs sont les suivantes :
	
	\begin{lstlisting}[caption={[] \\ \captionsource{\cite{mus:16}}}, label={lst:codeExample}, language=python, numbers=none]
		
		Flotteur = np.array([160, 140, 120, 100, 80, 60, 40, 20])  # en m
		
		c1=np.array([323, 343, 354, 361, 364, 365, 366, 366])
		c2=np.array([182, 233, 271, 299, 320, 346, 349, 357])
		c3=np.array([298, 323, 338, 348, 354, 358, 361, 363])
		c4=np.array([305, 328, 342, 352, 357, 360, 362, 364])
		c5=np.array([312, 333, 347, 355, 359, 361, 363, 364])
		c6=np.array([148, 209, 252, 284, 310, 330, 345, 355])
		c7=np.array([178, 230, 268, 297, 318, 334, 348, 355])
		c8=np.array([145, 204, 249, 281, 308, 328, 344, 353])
		c9=np.array([29, 195, 140, 176, 204, 225, 241, 252])
		
		
	\end{lstlisting}
		
		
		
% ----------------------------------------------

\chapter{Analyse des mesures}

% ############################################
\textbf{Tube de Venturi}

		% reponse a la question 4	
		\textcolor{blue}{Reponse 4-}
		
		%Demonstration...........................................
		
		\subsection*{Hypothèses}
		
		\begin{itemize}
			\item écoulement stationnaire
			\item fluide incompressible  
			\item pertes de charge négligeables
		\end{itemize}
		
		On applique la loi de Bernoulli entre les sections A et B :
		\[
		P_A + \frac{v_A^2}{2} +\rho g z_A = P_B + \frac{v_B^2}{2} +\rho g z_B
		\]
		
		\subsection*{Équation de continuité}
		
		Conservation de la masse :$	S_A v_A = S_B v_B  $
		
		
		
		\subsection*{Développement}
		
		
	
		En utilisant Bernoulli on a: 
		\[
		 P_A - P_B + \rho g (z_A - z_B) = \frac{1}{2} \rho \left( v_B^2 - v_A^2 \right)
		\]
		
		\[
			\rho g (h_A - h_B) = \frac{1}{2} \rho \left( v_B^2 - v_A^2 \right) = \frac{1}{2} \rho v_A^2 \left( \frac{S_A^2}{S_B^2} - 1 \right)
		\]
		
		\[
			 g (h_A - h_B) = \frac{1}{2} v_A^2 \left( \frac{d_A^4}{d_B^4} - 1 \right)
		\]
		
		\[
		v_A= \sqrt{ \frac{2 g (h_A h_B)  } {\left(   \frac{d_A ^4}{d_B^4} \left(  1 - \frac{d_B ^4}{d_A^4}  \right) \right)  } }
		\]
		
		\subsection*{Expression finale}
		
			On sait que
			
			 \[ Q = \frac{dvolume}{dt} = S* \frac{dx}{dt} = S_A v_A = \pi \frac{d_A^2}{4} v_A
			\]
		
		Donc on définit le coefficient :
		\[
		C = \frac{\pi d_B^2 \sqrt{g}}{\sqrt{8\left(1 - \left(\frac{d_B}{d_A}\right)^4\right)}}
		\]
		
		Donc :
		\[
		Q = C \sqrt{h_A - h_B}
		\]
		
		
		% reponse a la question 5	
		\textcolor{blue}{Reponse 5- Courbe expérimentale $Q = f\left( \sqrt{h_A -h_B} \right)$}
		
			\begin{figure}[H]
				\centering
				\includegraphics[width=0.5\textwidth]{question5.png}
				\caption{Trac de Q= A * Flotteur + B  }
				\label{fig:image3}
			\end{figure}
			
		\textbf{Observation:} On observe une répartition linéaire des points expérimentaux (points bleus). Cela indique que la relation entre le débit et la variable mesurée suit une loi de proportionnalité 
		
		
		
		
		% reponse a la question 6	
		\textcolor{blue}{Reponse 6-  Tracé pour chaque points des barres d'erreur }
		
	
		\begin{figure}[H]
			\centering
			\includegraphics[width=0.5\textwidth]{question6.png}
			\caption{Controle de la difference d'erreur  }
			\label{fig:image3}
		\end{figure}
		
		\textbf{Observation:} Il existe un écart systématique entre la théorie et l'expérience. La courbe théorique (orange) se situe nettement en dessous de nos points expérimentaux.\\
		
		\textbf{Commentaire sur l'écart:}  Dans la réalité, un fluide n'est pas parfait. Les frottements visqueux et les pertes de charge (que ce soit dans le Venturi ou le diaphragme) font que le débit réel diffère du débit théorique idéal \\
		On constate aussi que l'écart est dit significatif et ne peut pas être expliqué uniquement par l'imprécision de lecture.
		
		
		
		% reponse a la question 7	
		\textcolor{blue}{Reponse 7- Tracé sous la forme d'une courbe la prédiction théorique }
		
		\begin{figure}[H]
			\centering
			\includegraphics[width=0.5\textwidth]{question7.png}
			\caption{Courbe continue de la prédiction théorique  }
			\label{fig:image3}
		\end{figure}
		
		\textbf{Utilité : }  Elle sert de référence idéale. Elle montre que, selon Bernoulli, le débit devrait augmenter avec la racine carrée de la dénivelée piézométrique
		
		
		
		% reponse a la question 8	
		\textcolor{blue}{Reponse 8- Reprise du même travail pour B et C }
		
		Demonstration...........................................
		
		
		
		% reponse a la question 9	
		\textcolor{blue}{Reponse 9- Comparaison des résultats théoriques et expérimentaux }
		
		\begin{itemize}
			\item \textbf{Analyse qualitative :} 
			D'après les graphiques obtenues (notamment le tracé de $Q = f(\sqrt{h_A - h_B})$), on observe que les points expérimentaux suivent une tendance linéaire très marquée. Cela valide qualitativement la \textbf{loi de Bernoulli} pour un fluide incompressible, confirmant que le débit volumique $Q$ est bien proportionnel à la racine carrée de la dénivelée piézométrique : $Q \propto \sqrt{\Delta h}$ .
			
			\item \textbf{Analyse quantitative et écarts :} 
			Bien que la forme de la courbe soit respectée, on note un écart entre la courbe théorique et les points expérimentaux. Dans nos mesures, le débit expérimental est supérieur à la prédiction théorique simplifiée ($Q_{exp} > Q_{th}$). Cet écart peut s'expliquer par plusieurs facteurs :
			\begin{itemize}
				\item \textbf{Hypothèse du fluide parfait :} La théorie repose sur l'absence de viscosité. En réalité, l'eau est un fluide parfait. Les frottements visqueux aux parois et les turbulences dissipent de l'énergie sous forme de \textbf{pertes de charge} .
				\item \textbf{Coefficient de décharge ($C_c$) :} Pour compenser ces effets et la contraction du jet (vena contracta), il est nécessaire d'introduire un coefficient de correction $C_c$. La valeur théorique calculée avec la géométrie pure ne suffit pas à décrire parfaitement le système réel .
				\item \textbf{Incertitudes de mesure :} Les erreurs de lecture des hauteurs d'eau sur les tubes piézométriques (erreur de parallaxe) et l'imprécision lors du chronométrage manuel pour la mesure du débit au bac d'étalonnage contribuent à la dispersion des points .
			\end{itemize}
			
			\item \textbf{Conclusion :} 
			La loi de Bernoulli est vérifiée dans son principe de fonctionnement. Toutefois, pour une utilisation précise du tube de Venturi comme instrument de mesure industriel, un \textbf{étalonnage préalable} est indispensable afin de déterminer le coefficient de décharge réel du dispositif.
		\end{itemize}
% ############################################
\section{Diaphragme}

	% reponse a la question 10	
	\textcolor{blue}{Reponse 10: Evolution de profil de vitesse entre les section E et F  }\\	
		Entre E et G, les lignes de courant se resserrent au passage de l’orifice, puis s’écartent en aval. La vitesse augmente dans la zone contractée, puis diminue après l’orifice. 
		\\	
	% reponse a la question 11	
	\textcolor{blue}{Reponse 11: Explication de porquoi $ h_F < h_G $} \\
		Entre F et G, l’écoulement accélère au niveau de l’orifice, donc la pression statique diminue. Cela explique que la hauteur mesurée en F soit inférieure à celle en G.
	
	
	
	
	% reponse a la question 12	
	\textcolor{blue}{Reponse 12: : Loi théorique $Q = f(\sqrt{h_E - h_F})$}
		
		
		
		On applique le théorème de Bernoulli entre les sections E et F du diaphragme, en supposant un écoulement stationnaire, incompressible et sans pertes de charge :
		
		\begin{equation}
			\frac{P_E}{\rho g} + \frac{v_E^2}{2g} + z_E = \frac{P_F}{\rho g} + \frac{v_F^2}{2g} + z_F
		\end{equation}
		
		Si les prises sont au même niveau ($z_E = z_F$), cela donne :
		\begin{equation}
			P_E - P_F = \frac{\rho}{2} (v_F^2 - v_E^2)
		\end{equation}
		
		Par conservation de la masse (équation de continuité) :
		\begin{equation}
			S_E v_E = S_F v_F \quad \Rightarrow \quad v_E = \frac{Q}{S_E}, \quad v_F = \frac{Q}{S_F}
		\end{equation}
		
		En substituant :
		\begin{equation}
			\rho g (h_E - h_F) = \frac{\rho}{2} Q^2 \left( \frac{1}{S_F^2} - \frac{1}{S_E^2} \right)
		\end{equation}
		
		Donc :
		\begin{equation}
			Q = \frac{S_F}{\sqrt{1 - \left(\frac{S_F}{S_E}\right)^2}} \sqrt{2g (h_E - h_F)}
		\end{equation}
		
		
		
	% reponse a la question 13	
	\textcolor{blue}{Reponse 13:  Perte de charge du diaphragme }
	
		La perte de charge entre E et G s'obtient par différence hydrostatique :
		\begin{equation}
			\Delta X = \rho g (h_E - h_G)
		\end{equation}
		
		Cette perte correspond à l'énergie mécanique dissipée par les turbulences et les frottements.
	
	
	
	
	
	% reponse a la question 14	
	\textcolor{blue}{ Reponse 14:  La perte de charge est-elle négligeable ?}
	
	
	On compare le rapport :
	\begin{equation}
		\frac{h_E - h_G}{h_E - h_F}
	\end{equation}
	
	Si ce rapport est $\ll 1$, la perte est négligeable. En pratique, pour un diaphragme, ce rapport est souvent de l'ordre de 10-30\%, donc non négligeable.
	
	
	
	% reponse a la question 15	
	\textcolor{blue}{Reponse 15: Détermination du coefficient $C_c$ }
	
			\begin{figure}[H]
			\centering
			\includegraphics[width=0.5\textwidth]{question14.png}
			\caption{ Tracé de la courbe expérimentale  }
			\label{fig:image3}
		\end{figure}
		
		Le facteur de correction donne:  $ C_c =  0.12383292187219626 $
		
		La loi corrigée est :
		\begin{equation}
			Q = C_c \frac{\pi d_F^2}{\sqrt{8}} \sqrt{\frac{g (h_E - h_F)}{1 - \left(\frac{d_F}{d_E}\right)^4}}
		\end{equation}
		
		Le coefficient $C_c$ se détermine expérimentalement par :
		\begin{equation}
			C_c = \frac{Q_{\text{exp}} \sqrt{8} \sqrt{1 - \left(\frac{d_F}{d_E}\right)^4}}{\pi d_F^2 \sqrt{g (h_E - h_F)}}
		\end{equation}
		

	% ############################################
	\section{Rotametre}
	
		% reponse a la question 16	
		\textbf{Reponses 16 : Lignes de courant dans le rotamètre  }
		Les lignes de courant se resserrent dans l'espace annulaire entre le flotteur et la paroi conique du rotamètre. La vitesse est maximale dans cette zone de passage réduit.
		---------------------- Images -------------------------------------- \\
		% reponse a la question 17
	\textcolor{blue}{Reponse 17: Forces agissant sur le flotteur  }
	
	À l'équilibre, le flotteur subit trois forces principales :
	
	\begin{enumerate}
		\item \textbf{Poids} : $P =- m g$ (vers le bas)
		\item \textbf{Poussée d'Archimède} : $P_A = \rho_f V_{\text{immergé}} g$ (vers le haut)
		\item \textbf{Force hydrodynamique} : due à la différence de pression autour du flotteur (vers le haut)
	\end{enumerate}
	
	Équilibre : $P = P_A + F_{\text{hyd}}$
	
	
	% reponse a la question 19	
	\textcolor{blue}{Reponse 18 :  Augmentation du débit à $L$ constant}
	
		Si on augmente $Q$ en maintenant $L$ constant :
		\begin{itemize}
			\item La vitesse dans l'espace annulaire augmente
			\item La différence de pression $\Delta P$ autour du flotteur augmente
			\item $F_{\text{hyd}}$ augmente
			\item Le flotteur n'est plus en équilibre et tend à monter
		\end{itemize}
	
	
	% reponse a la question 19	
	\textcolor{blue}{Reponse 19: Pourquoi le flotteur monte-t-il ? }
	
	
	Quand on relâche le flotteur :
	\begin{enumerate}
		\item $F_{\text{hyd}} > P - P_A$ $\Rightarrow$ le flotteur monte
		\item En montant, la section annulaire $S_{\text{ann}}$ augmente
		\item La vitesse $v = Q/S_{\text{ann}}$ diminue
		\item $\Delta P$ et $F_{\text{hyd}}$ diminuent jusqu'au nouvel équilibre
	\end{enumerate}
	
	
	% reponse a la question 20	
	\textcolor{blue}{Reponse 20:  Pertes de charge du rotamètre}
	
	Les pertes de charge sont :
	\begin{equation}
		\Delta X = \rho g (h_H - h_I)
	\end{equation}
	
	Elles sont dues à :
	\begin{itemize}
		\item Frottements visqueux dans l'espace annulaire
		\item Turbulences autour du flotteur
		\item Changements brusques de section
	\end{itemize}
	
	% reponse a la question 21	
	\textcolor{blue}{Reponse 21: Courbe $\Delta X = f(Q)$ }
	
	\begin{figure}[H]
		\centering
		\includegraphics[width=0.5\textwidth]{question21.png}
		\caption{Tracé de la courbe expérimentale $\Delta X = f(Q)$}
		\label{fig:deltaX_Q}
	\end{figure}
	
		
		
		\begin{itemize}
			\item \textbf{Observations :}
			La courbe expérimentale représentant l'évolution de la perte de charge $\Delta X$ (définie par $\rho g(h_H - h_I)$) en fonction du débit volumique $Q$ présente une allure \textbf{parabolique}. Contrairement aux courbes d'étalonnage précédentes, la pente de ce tracé n'est pas constante et croît avec le débit.
			
			\item \textbf{Interprétation physique :}
			Cette forme quadratique est caractéristique des \textbf{pertes de charge singulières} en régime inertiel. Pour un organe de mesure comme le rotamètre, le fluide subit un rétrécissement puis un élargissement brusque en contournant le flotteur. La théorie prévoit que la perte d'énergie mécanique suit une loi proportionnelle au carré de la vitesse :
			\[ \Delta X = \frac{1}{2} \rho u^2 K \]
			Sachant que la vitesse $u$ est proportionnelle au débit volumique ($Q = u \cdot S$), on obtient bien une relation de type $\Delta X = \alpha Q^2$.
			
			\item \textbf{Origine du phénomène :}
			La dissipation d'énergie ne provient pas ici du frottement visqueux linéaire, mais du brassage intense du fluide et de la création de \textbf{tourbillons} dans le sillage du flotteur. Ces phénomènes turbulents dissipent une énergie mécanique proportionnelle à l'énergie cinétique du jet annulaire, expliquant ainsi la dépendance au carré du débit.
			
			\item \textbf{Conclusion :}
			Le résultat obtenu valide le caractère inertiel de l'écoulement au sein du rotamètre. La perte de charge n'est pas négligeable et augmente de façon significative pour les hauts débits.
		\end{itemize}
		
		
	
		% reponse a la question 22	
	\textcolor{blue}{Reponse 22: Courbe $\Delta X = f(Q)$ }
	
	---------------------- Images -------------------------------------- \\
	
	On ajuste expérimentalement une loi du type :
	\begin{equation}
		Q = a L + b
	\end{equation}
	
	ou plus généralement :
	\begin{equation}
		Q = k L^n
	\end{equation}
	
	où $a$, $b$, $k$, $n$ sont déterminés par régression linéaire ou non linéaire des mesures.
	

		\end{document}
	
\pagebreak

% ############################################
\section{Lists}

Unordered Lists:
\begin{itemize}
\item This is an unordered list. 
\item Item 2.
\item It has three items.
\end{itemize}

Ordered List:
\begin{enumerate}
\item This is an ordered list.
\item Item 2.
\item It has three items.
\end{enumerate}

Ordered List (alphabetical):
\begin{enumerate}[label=\Alph*.]
\item This is an ordered list.
\item Item 2.
\item It has three items.
\end{enumerate}

% ----------------------------------------------
\chapter{Figures}\label{ch:figures}

% ############################################
\section{Images}\label{sec:images}

\autoref{fig:image1} shows how to display images.

\begin{figure}[H]
	\centering
%	\includegraphics[width=0.5\textwidth]{f1.jpg}
    \caption{Image}

	\label{fig:image1}
\end{figure}

\begin{figure}[H]
	\centering
%	\includegraphics[width=0.5\textwidth]{f4.jpg}
    \caption{Image with Source}
    \captionsource{\cite{mus:16}}	
	\label{fig:image2}
\end{figure}

\begin{figure}[H]
	\centering
%	\includegraphics[width=0.5\textwidth]{f7.jpg}
    \caption{Image with Source and Link}
    \captionsource[https://example.org]{J. Doe}	
	\label{fig:image3}
\end{figure}

% ############################################



















% ----------------------------------------------

\chapter{Conclusion}\label{ch:conclusion}

% ############################################
\section{Mathematics}\label{sec:maths}

\subsection{Definitions}

\begin{define}	
A function $f: X \rightarrow Y$ is injective if and only if for all $x_1, x_2 \in X$, $x_1 \neq x_2$ implies $f(x_1) \neq f(x_2)$.
\end{define}

\subsection{Formulas}

Formulas can be included inline, e.g. $K_E^{CID}$ to $S$ or as an own block.
$$ otp = f_k(i, k), \qquad i = 1, ..., n $$

\pagebreak
% ############################################
\section{Code Listings}\label{sec:code}

Code can be displayed inline, e.g. \code{UPDATE} or \code{factorial(n)}.

Code listings can have numbered lines, captions, syntax highlighting, ... .

\begin{lstlisting}[caption={[Code Without Line Numbers] Code Without Line Numbers \\ \captionsource{\cite{mus:16}}}, label={lst:codeExample}, language=python, numbers=none]
 0  1  2  3  4  5  6  7  8  9 10 11 12 13 14 15 16 17 
 A  B  C  D  E  F  G  H  I  J  K  L  M  N  O  P  Q  R  
\end{lstlisting}

\begin{lstlisting}[caption={Code With Line Numbers}, label={lst:factorial}, language=Java]
public int factorial(int n) {
	if (n == 1) {
		return 1;
	} else {
		return n * factorial(n - 1) ;
	}
}
\end{lstlisting}







